\documentclass[a4paper, 14pt]{extarticle}

% Поля
%--------------------------------------
\usepackage{geometry}
\geometry{a4paper,tmargin=2cm,bmargin=2cm,lmargin=3cm,rmargin=1cm}
%--------------------------------------

%Russian-specific packages
%--------------------------------------
\usepackage[T2A]{fontenc}
\usepackage[utf8]{inputenc}
\usepackage[english, main=russian]{babel}
%--------------------------------------

\usepackage{textcomp}

% Красная строка
%--------------------------------------
\usepackage{indentfirst}
%--------------------------------------

%Graphics
%--------------------------------------
\usepackage{graphicx}
\graphicspath{ {./images/} }
\usepackage{wrapfig}
%--------------------------------------

% Полуторный интервал
%--------------------------------------
\usepackage{setspace}
\linespread{1.3}
%--------------------------------------

%Выравнивание и переносы
%--------------------------------------
% Избавляемся от переполнений
\sloppy
% Запрещаем разрыв страницы после первой строки абзаца
\clubpenalty=10000
% Запрещаем разрыв страницы после последней строки абзаца
\widowpenalty=10000
%--------------------------------------

%Списки
\usepackage{enumitem}

%Подписи
\usepackage{caption}

%Таблицы и графики
\usepackage{float}
\usepackage{array}
\usepackage{tikz}
\usepackage{pgfplots}
\pgfplotsset{compat=1.18}

%Листинги
\usepackage{listings}
\usepackage{xcolor}

\lstdefinelanguage{Go}{
    morekeywords={break,default,func,interface,select,case,defer,go,map,struct,chan,else,goto,package,switch,const,fallthrough,if,range,type,continue,for,import,return,var},
    sensitive=true,
    morecomment=[l]{//},
    morecomment=[s]{/*}{*/},
    morestring=[b]"
}

\lstset{
    language=Go,
    basicstyle=\ttfamily\small,
    breaklines=true,
    columns=fullflexible,
    frame=single,
    inputencoding=utf8,
    extendedchars=true,
    keywordstyle=\bfseries,
    numbers=left,
    numbersep=5pt,
    xleftmargin=25pt,
    xrightmargin=25pt,
    showstringspaces=false,
    literate=
    {А}{{\CYRA}}1 {Б}{{\CYRB}}1 {В}{{\CYRV}}1 {Г}{{\CYRG}}1 {Д}{{\CYRD}}1
    {Е}{{\CYRE}}1 {Ё}{{\CYRYO}}1 {Ж}{{\CYRZH}}1 {З}{{\CYRZ}}1 {И}{{\CYRI}}1
    {Й}{{\CYRISHRT}}1 {К}{{\CYRK}}1 {Л}{{\CYRL}}1 {М}{{\CYRM}}1 {Н}{{\CYRN}}1
    {О}{{\CYRO}}1 {П}{{\CYRP}}1 {Р}{{\CYRR}}1 {С}{{\CYRS}}1 {Т}{{\CYRT}}1
    {У}{{\CYRU}}1 {Ф}{{\CYRF}}1 {Х}{{\CYRH}}1 {Ц}{{\CYRC}}1 {Ч}{{\CYRCH}}1
    {Ш}{{\CYRSH}}1 {Щ}{{\CYRSHCH}}1 {Ъ}{{\CYRHRDSN}}1 {Ы}{{\CYRERY}}1
    {Ь}{{\CYRSFTSN}}1 {Э}{{\CYREREV}}1 {Ю}{{\CYRYU}}1 {Я}{{\CYRYA}}1
    {а}{{\cyra}}1 {б}{{\cyrb}}1 {в}{{\cyrv}}1 {г}{{\cyrg}}1 {д}{{\cyrd}}1
    {е}{{\cyre}}1 {ё}{{\cyryo}}1 {ж}{{\cyrzh}}1 {з}{{\cyrz}}1 {и}{{\cyri}}1
    {й}{{\cyrishrt}}1 {к}{{\cyrk}}1 {л}{{\cyrl}}1 {м}{{\cyrm}}1 {н}{{\cyrn}}1
    {о}{{\cyro}}1 {п}{{\cyrp}}1 {р}{{\cyrr}}1 {с}{{\cyrs}}1 {т}{{\cyrt}}1
    {у}{{\cyru}}1 {ф}{{\cyrf}}1 {х}{{\cyrh}}1 {ц}{{\cyrc}}1 {ч}{{\cyrch}}1
    {ш}{{\cyrsh}}1 {щ}{{\cyrshch}}1 {ъ}{{\cyrhrdsn}}1 {ы}{{\cyrery}}1
    {ь}{{\cyrsftsn}}1 {э}{{\cyrerev}}1 {ю}{{\cyryu}}1 {я}{{\cyrya}}1
}

%Рисунки
%--------------------------------------
\DeclareCaptionLabelSeparator*{emdash}{~--- }
\captionsetup[figure]{labelsep=emdash,font=onehalfspacing,position=bottom}
%--------------------------------------

\usepackage{tempora}


%--------------------------------------
%   НАЧАЛО ДОКУМЕНТА
%--------------------------------------

\begin{document}

%--------------------------------------
%   ТИТУЛЬНЫЙ ЛИСТ
%--------------------------------------
\begin{titlepage}
\thispagestyle{empty}
\newpage


%Шапка титульного листа
%--------------------------------------
\vspace*{-60pt}
\hspace{-65pt}
\begin{minipage}{0.3\textwidth}
\hspace*{-20pt}\centering
\includegraphics[width=\textwidth]{emblem}
\end{minipage}
\begin{minipage}{0.67\textwidth}\small \textbf{
\vspace*{-0.7ex}
\hspace*{-6pt}\centerline{Министерство науки и высшего образования Российской Федерации}
\vspace*{-0.7ex}
\centerline{Федеральное государственное бюджетное образовательное учреждение }
\vspace*{-0.7ex}
\centerline{высшего образования}
\vspace*{-0.7ex}
\centerline{<<Московский государственный технический университет}
\vspace*{-0.7ex}
\centerline{имени Н.Э. Баумана}
\vspace*{-0.7ex}
\centerline{(национальный исследовательский университет)>>}
\vspace*{-0.7ex}
\centerline{(МГТУ им. Н.Э. Баумана)}}
\end{minipage}
%--------------------------------------

%Полосы
%--------------------------------------
\vspace{-25pt}
\hspace{-35pt}\rule{\textwidth}{2.3pt}

\vspace*{-20.3pt}
\hspace{-35pt}\rule{\textwidth}{0.4pt}
%--------------------------------------

\vspace{1.5ex}
\hspace{-35pt} \noindent \small ФАКУЛЬТЕТ\hspace{80pt} <<Информатика и системы управления>>

\vspace*{-16pt}
\hspace{47pt}\rule{0.83\textwidth}{0.4pt}

\vspace{0.5ex}
\hspace{-35pt} \noindent \small КАФЕДРА\hspace{50pt} <<Теоретическая информатика и компьютерные технологии>>

\vspace*{-16pt}
\hspace{30pt}\rule{0.866\textwidth}{0.4pt}

\vspace{11em}

\begin{center}
\Large {\bf Лабораторная работа № 1} \\

\large {\bf по курсу <<Разработка параллельных и распределённых программ>>} \\[0.8cm]

\large <<Распараллеливание алгоритма вычисления
произведения двух матриц>>
\end{center}
\normalsize

\vspace{8em}\begin{flushright}
  {Студент группы ИУ9-51Б Махова А. А. \hspace*{15pt}\\
  \vspace{2ex}
  Преподаватель Царев А. С.\hspace*{15pt}}
\end{flushright}

\bigskip

\vfill


\begin{center}
\textsl{Москва 2026}
\end{center}
\end{titlepage}


% =================================================
% ОГЛАВЛЕНИЕ
% =================================================

\setcounter{page}{2}

\tableofcontents

\newpage


% =================================================
% 1. ПОСТАНОВКА ЗАДАЧИ
% =================================================

\section{Постановка задачи}

Целью лабораторной работы является исследование последовательного
и параллельного алгоритмов умножения квадратных матриц,
а также сравнение времени их выполнения.

Даны две квадратные матрицы $A$ и $B$ размерности $n \times n$.
Необходимо получить результирующую матрицу $C$ той же размерности.

Элемент результирующей матрицы вычисляется по формуле:

\[
c_{ij}
=
\sum_{k=1}^{n}
a_{ik}b_{kj}.
\]

В ходе выполнения лабораторной работы необходимо:

\begin{enumerate}[
    label=\arabic*.,
    itemsep=0pt,
    topsep=3pt,
    parsep=0pt,
    partopsep=0pt
]

    \item Реализовать стандартный последовательный алгоритм
    умножения двух квадратных матриц.

    \item Выполнить вычисление элементов результирующей матрицы
    последовательно по строкам.

    \item Выполнить вычисление элементов результирующей матрицы
    последовательно по столбцам.

    \item Сравнить время работы последовательного вычисления
    по строкам и по столбцам.

    \item Разделить результирующую матрицу на примерно равные
    прямоугольные подматрицы.

    \item Реализовать параллельное вычисление элементов
    результирующей матрицы.

    \item Выполнить вычисления для различных вариантов разбиения
    матрицы и различного количества рабочих потоков.

    \item Организовать ожидание завершения всех параллельных
    вычислений.

    \item После окончания вычислений сравнить полученный результат
    с результатом стандартного последовательного алгоритма.

    \item Измерить время выполнения различных вариантов алгоритма
    и сравнить полученные результаты.

\end{enumerate}

Использование библиотечных функций для непосредственного
умножения матриц не предусматривается.


% =================================================
% 2. ЛИСТИНГ ПРОГРАММЫ
% =================================================

\section{Листинг программы}

Ниже приведён полный листинг программы на языке Go,
использованной для проведения вычислений.

\begin{lstlisting}
package main

import (
	"flag"
	"fmt"
	"math"
	"math/rand"
	"runtime"
	"sync"
	"time"
)

type Matrix struct {
	n    int
	data []float64
}

func NewMatrix(n int) *Matrix {
	return &Matrix{n: n, data: make([]float64, n*n)}
}

func (m *Matrix) At(i, j int) float64 {
	return m.data[i*m.n+j]
}

func (m *Matrix) Set(i, j int, v float64) {
	m.data[i*m.n+j] = v
}

func (m *Matrix) FillRand(r *rand.Rand) {
	for i := range m.data {
		m.data[i] = r.Float64()*2 - 1
	}
}

// 1. Последовательное умножение: вычисляем C по строкам.
func MulRows(A, B *Matrix) *Matrix {
	n := A.n
	C := NewMatrix(n)

	for i := 0; i < n; i++ {
		for j := 0; j < n; j++ {
			sum := 0.0
			for k := 0; k < n; k++ {
				sum += A.At(i, k) * B.At(k, j)
			}
			C.Set(i, j, sum)
		}
	}
	return C
}

// 2. То же самое произведение, но элементы C вычисляются по столбцам.
func MulColumns(A, B *Matrix) *Matrix {
	n := A.n
	C := NewMatrix(n)

	for j := 0; j < n; j++ {
		for i := 0; i < n; i++ {
			sum := 0.0
			for k := 0; k < n; k++ {
				sum += A.At(i, k) * B.At(k, j)
			}
			C.Set(i, j, sum)
		}
	}
	return C
}

// Вычисление одной прямоугольной подматрицы C.
func multiplyBlock(A, B, C *Matrix, row0, row1, col0, col1 int) {
	n := A.n
	for i := row0; i < row1; i++ {
		for j := col0; j < col1; j++ {
			sum := 0.0
			for k := 0; k < n; k++ {
				sum += A.At(i, k) * B.At(k, j)
			}
			C.Set(i, j, sum)
		}
	}
}

// 3. Параллельное умножение.
// Матрица C делится на rowParts x colParts прямоугольных блоков.
// Каждый блок вычисляется своей goroutine.
func MulParallelGrid(A, B *Matrix, rowParts, colParts int) *Matrix {
	n := A.n
	C := NewMatrix(n)

	if rowParts < 1 {
		rowParts = 1
	}
	if colParts < 1 {
		colParts = 1
	}
	if rowParts > n {
		rowParts = n
	}
	if colParts > n {
		colParts = n
	}

	var wg sync.WaitGroup

	for br := 0; br < rowParts; br++ {
		row0 := br * n / rowParts
		row1 := (br + 1) * n / rowParts

		for bc := 0; bc < colParts; bc++ {
			col0 := bc * n / colParts
			col1 := (bc + 1) * n / colParts

			wg.Add(1)
			go func(r0, r1, c0, c1 int) {
				defer wg.Done()
				multiplyBlock(A, B, C, r0, r1, c0, c1)
			}(row0, row1, col0, col1)
		}
	}

	// Основная goroutine ждёт завершения всех рабочих goroutine.
	wg.Wait()
	return C
}

func maxAbsDiff(X, Y *Matrix) float64 {
	maxDiff := 0.0
	for i := range X.data {
		d := math.Abs(X.data[i] - Y.data[i])
		if d > maxDiff {
			maxDiff = d
		}
	}
	return maxDiff
}

func speedup(base, cur time.Duration) float64 {
	if cur <= 0 {
		return math.Inf(1)
	}
	return float64(base) / float64(cur)
}

func main() {
	n := flag.Int("n", 500, "размер квадратной матрицы n x n")
	flag.Parse()

	fmt.Printf("GOMAXPROCS = %d\n", runtime.GOMAXPROCS(0))
	fmt.Printf("Matrix size: %d x %d\n\n", *n, *n)

	r := rand.New(rand.NewSource(1))

	A := NewMatrix(*n)
	B := NewMatrix(*n)
	A.FillRand(r)
	B.FillRand(r)

	// Последовательно по строкам.
	start := time.Now()
	Crows := MulRows(A, B)
	rowsTime := time.Since(start)
	fmt.Printf("[SEQ rows]    time = %v\n", rowsTime)

	// Последовательно по столбцам.
	start = time.Now()
	Ccols := MulColumns(A, B)
	colsTime := time.Since(start)
	fmt.Printf("[SEQ columns] time = %v   max|diff| = %.3e\n\n",
		colsTime, maxAbsDiff(Crows, Ccols))

	// Разные прямоугольные разбиения C.
	tests := []struct {
		rows int
		cols int
	}{
		{1, 2}, // 2 goroutines
		{2, 2}, // 4 goroutines
		{2, 4}, // 8 goroutines
		{4, 4}, // 16 goroutines
		{4, 8}, // 32 goroutines
		{8, 8}, // 64 goroutines
	}

	for _, t := range tests {
		workers := t.rows * t.cols

		start = time.Now()
		Cpar := MulParallelGrid(A, B, t.rows, t.cols)
		parTime := time.Since(start)

		diff := maxAbsDiff(Crows, Cpar)

		fmt.Printf(
			"[PAR %2dx%-2d = %2d workers] time = %-12v speedup = %6.2fx   max|diff| = %.3e\n",
			t.rows,
			t.cols,
			workers,
			parTime,
			speedup(rowsTime, parTime),
			diff,
		)
	}

	fmt.Println("\nЕсли max|diff| близко к 0, результат вычислен правильно.")
	fmt.Println("В Go нет переносимого API для задания приоритета отдельной goroutine,")
	fmt.Println("поэтому эксперимент с high/low priority для goroutine обычно не выполняется.")
}
\end{lstlisting}


% =================================================
% 3. РЕЗУЛЬТАТЫ ВЫЧИСЛЕНИЙ
% =================================================

\section{Результаты вычислений}

Для исследования производительности были проведены
последовательные и параллельные вычисления произведения
двух квадратных матриц. Размер матриц задавался параметром
\texttt{-n}; в программе по умолчанию используется значение
$n = 500$.

Сначала выполнялось последовательное умножение по строкам,
результат которого использовался в качестве эталонного.
После этого выполнялось последовательное вычисление по столбцам,
а затем параллельные вычисления с различными вариантами
разбиения результирующей матрицы.

\begin{table}[H]

\centering

\caption{Время выполнения алгоритмов}

\resizebox{\textwidth}{!}{
\begin{tabular}{|c|c|c|c|c|}
\hline

\textbf{Режим}
&
\textbf{Разбиение}
&
\textbf{Число рабочих}
&
\textbf{Время, с}
&
\textbf{Ускорение}
\\
\hline

Последовательный по строкам
&
---
&
1
&
\textbf{ВСТАВИТЬ}
&
1
\\
\hline

Последовательный по столбцам
&
---
&
1
&
\textbf{ВСТАВИТЬ}
&
---
\\
\hline

Параллельный
&
$1 \times 2$
&
2
&
\textbf{ВСТАВИТЬ}
&
\textbf{ВСТАВИТЬ}
\\
\hline

Параллельный
&
$2 \times 2$
&
4
&
\textbf{ВСТАВИТЬ}
&
\textbf{ВСТАВИТЬ}
\\
\hline

Параллельный
&
$2 \times 4$
&
8
&
\textbf{ВСТАВИТЬ}
&
\textbf{ВСТАВИТЬ}
\\
\hline

Параллельный
&
$4 \times 4$
&
16
&
\textbf{ВСТАВИТЬ}
&
\textbf{ВСТАВИТЬ}
\\
\hline

Параллельный
&
$4 \times 8$
&
32
&
\textbf{ВСТАВИТЬ}
&
\textbf{ВСТАВИТЬ}
\\
\hline

Параллельный
&
$8 \times 8$
&
64
&
\textbf{ВСТАВИТЬ}
&
\textbf{ВСТАВИТЬ}
\\
\hline

\end{tabular}
}

\end{table}

Во всех рассмотренных вариантах после окончания вычислений
производилась проверка результата относительно
последовательного алгоритма. Если значение максимального
абсолютного отклонения близко к нулю, результат считается
полученным корректно.


% =================================================
% 4. ГРАФИК
% =================================================

\section{График зависимости времени от числа рабочих потоков}

Для анализа производительности был построен график зависимости
времени параллельного вычисления от количества рабочих goroutine.

По оси абсцисс откладывается количество рабочих goroutine,
а по оси ординат --- время выполнения алгоритма в секундах.

\begin{figure}[H]

\centering

\begin{tikzpicture}

\begin{axis}[
    width=0.90\textwidth,
    height=9cm,
    xlabel={Количество рабочих goroutine},
    ylabel={Время вычисления, с},
    grid=major,
    xtick={2,4,8,16,32,64},
    ymin=0,
    legend pos=north east
]

\addplot[
    mark=*,
    thick
]
coordinates {

    % Заменить нули на полученные результаты измерений.

    (2,0)
    (4,0)
    (8,0)
    (16,0)
    (32,0)
    (64,0)

};

\legend{Параллельное вычисление}

\end{axis}

\end{tikzpicture}

\caption{Зависимость времени вычисления от количества рабочих goroutine}

\end{figure}


% =================================================
% 5. СРАВНЕНИЕ
% =================================================

\section{Сравнение последовательного и параллельного вариантов}

В последовательном варианте все элементы результирующей матрицы
вычисляются одной goroutine.

При последовательном вычислении по строкам сначала полностью
вычисляются элементы одной строки результирующей матрицы,
после чего программа переходит к следующей строке.

При последовательном вычислении по столбцам сначала вычисляются
элементы одного столбца, после чего выполняется переход к
следующему столбцу.

Количество арифметических операций в обоих случаях одинаково,
однако изменяется порядок обращения к данным в памяти,
что может влиять на время выполнения.

В параллельном варианте результирующая матрица разделяется
на прямоугольные блоки. Каждый блок вычисляется отдельной
goroutine. Для ожидания завершения всех goroutine используется
\texttt{sync.WaitGroup}.

Это позволяет выполнять несколько независимых частей вычисления
одновременно. При увеличении количества goroutine время выполнения
может уменьшаться за счёт использования нескольких ядер процессора.

При этом ускорение не обязано быть прямо пропорционально количеству
goroutine, поскольку на время работы также влияют накладные расходы
на создание и планирование goroutine, синхронизацию, работу с памятью
и количество реально доступных процессорных ядер.

В языке Go отсутствует переносимый API для задания приоритета
отдельной goroutine, поэтому эксперимент с высоким и низким
приоритетом goroutine в данной программе не выполнялся.

Конкретные выводы об эффективности различных разбиений делаются
на основании результатов измерений, представленных в таблице
и на графике.


% =================================================
% 6. ХАРАКТЕРИСТИКИ УСТРОЙСТВА
% =================================================

\section{Характеристики устройства}

Экспериментальные вычисления проводились на устройстве
со следующими характеристиками:

\begin{itemize}

    \item Операционная система ---
    macOS 15.2;

    \item Процессор ---
    Apple M4;

    \item Количество ядер процессора ---
    10;

    \item Объём оперативной памяти ---
    16 ГБ;

    \item Архитектура процессора ---
    arm64;

    \item Версия языка Go ---
    \textbf{ВСТАВИТЬ};

    \item Значение \texttt{GOMAXPROCS} ---
    10.

\end{itemize}

Перед началом вычислений программа самостоятельно выводит
часть характеристик системы:

\begin{lstlisting}
fmt.Printf("GOMAXPROCS = %d\n", runtime.GOMAXPROCS(0))
fmt.Printf("Matrix size: %d x %d\n\n", *n, *n)
\end{lstlisting}


% =================================================
% 7. ВЫВОД
% =================================================

\section{Вывод}

В ходе выполнения лабораторной работы был реализован алгоритм
умножения квадратных матриц на языке Go.

Были рассмотрены два последовательных варианта: вычисление элементов
результирующей матрицы по строкам и вычисление по столбцам.

Для реализации параллельного алгоритма результирующая матрица
разделялась на примерно равные прямоугольные блоки. Для вычисления
каждого блока создавалась отдельная goroutine.

Синхронизация параллельных вычислений была реализована при помощи
механизма \texttt{sync.WaitGroup}. Основная goroutine ожидала
окончания работы всех созданных рабочих goroutine.

Корректность результатов проверялась путём сравнения каждой
полученной матрицы с эталонной матрицей, вычисленной последовательным
алгоритмом по строкам.

Были проведены измерения времени работы программы для различных
вариантов разбиения результирующей матрицы. Для параллельных
вариантов также рассчитывалось ускорение относительно
последовательного вычисления по строкам.

Полученные результаты позволяют сравнить последовательный и
параллельный способы вычисления произведения матриц и определить
влияние количества рабочих goroutine на время выполнения программы.

\end{document}
